\section{Experiments}
\label{sec:experiments}

% Setup v1 (2026-06-10). Sous-sections independantes des resultats finaux :
% Datasets, Baselines, Metrics, Evaluation Protocol, Implementation Details.
% Les resultats / ablations restent en placeholder (attendent les chiffres).
% Materiel ANONYMISE (double-blind) : "consumer-grade GPU, 6 GB VRAM".

\subsection{Datasets}
\label{sec:exp:datasets}
We evaluate on \textbf{13 evaluation sets} that span news, social media, spoken queries, encyclopedic text, biomedical abstracts and manufacturing literature, following the dataset selection of Genest et al. \cite{genest2025owner}.
CrossNER \cite{liu2021crossner} contributes five domain-specific sets, evaluated separately (artificial intelligence, literature, music, politics, science).
The remaining eight are CoNLL-2003 \cite{tjongkimsang2003conll}, WNUT-17 \cite{derczynski2017wnut}, the MIT-Restaurant and MIT-Movie spoken-query corpora \cite{liu2013asgard}, FabNER \cite{kumar2022fabner}, the BioNLP/JNLPBA-2004 biomedical task \cite{kim2004jnlpba}, and the GUM \cite{zeldes2017gum} and GENTLE \cite{aoyama2023gentle} multi-genre corpora.
Label granularity ranges from four coarse types in CoNLL-2003 to twelve fine-grained categories in FabNER, including opaque domain acronyms.
Two datasets used by \cite{genest2025owner}, GENIA \cite{kim2003genia} and i2b2, are license-gated and therefore not re-measured here; BioNLP/JNLPBA-2004 serves as our biomedical proxy.
To homogenise evaluation cost across systems, we cap each test split at $1000$ sentences (the full split is used when it is smaller).
\autoref{tab:datasets} summarises the benchmark.

\begin{table}[t]
\centering
\caption{The 13 evaluation sets. ``Test'' is the number of test sentences used (capped at 1000).}
\label{tab:datasets}
\begin{tabular}{lllr}
\toprule
Dataset & Abbr. & Domain & Test \\
\midrule
CrossNER-AI         & AI    & Encyclopedic (AI)         & 430 \\
CrossNER-Literature & Lit   & Encyclopedic (literature) & 416 \\
CrossNER-Music      & Mus   & Encyclopedic (music)      & 465 \\
CrossNER-Politics   & Pol   & Encyclopedic (politics)   & 651 \\
CrossNER-Science    & Sci   & Encyclopedic (science)    & 543 \\
WNUT-17             & WNUT  & Social media              & 689 \\
MIT-Restaurant      & Rest  & Spoken queries            & 1000 \\
MIT-Movie           & Movie & Spoken queries            & 1000 \\
FabNER              & Fab   & Manufacturing             & 1000 \\
BioNLP/JNLPBA-2004  & Bio   & Biomedical                & 1000 \\
CoNLL-2003          & CoNLL & News                      & 1000 \\
GUM                 & GUM   & Multi-genre               & 1000 \\
GENTLE              & GEN   & Multi-genre (OOD)         & 1000 \\
\bottomrule
\end{tabular}
\end{table}

\subsection{Baselines}
\label{sec:exp:baselines}
We compare Opener against representative open-world NER systems run end-to-end on the same splits.
GLiNER \cite{zaratiana2024gliner} is evaluated in its small, medium and large variants to trace a size/efficiency trade-off.
GNER \cite{ding2024gner} is a generative tagger built on T5 \cite{raffel2020t5}; we run the T5-base checkpoint.
We further compare against OWNER \cite{genest2025owner}, executed in its original environment.
Instruction-tuned and few-shot large-language-model NER systems, namely UniversalNER \cite{zhou2024universalner}, GoLLIE \cite{sainz2024gollie} and ChatIE \cite{wei2024chatie}, rely on 7B-parameter backbones \cite{touvron2023llama} that do not fit our hardware budget even under 4-bit NF4 quantization \cite{dettmers2023qlora}; for these we report the figures published by their authors rather than re-measuring them.

\subsection{Metrics}
\label{sec:exp:metrics}
Our primary metric is the Adjusted Mutual Information (AMI) \cite{vinh2010ami} between predicted and gold type assignments, following \cite{genest2025owner}, complemented by macro-averaged $F_1$ and accuracy.
Because open-world predictions need not share the gold label vocabulary, AMI is computed over the aligned set of mentions, which makes it well suited to comparing dynamically named types.
Beyond quality, we report two efficiency axes: inference \emph{speed}, as the median (p50) latency per sentence, and \emph{energy}, as kilowatt-hours and grams of CO\textsubscript{2}eq per run.

\subsection{Evaluation Protocol}
\label{sec:exp:protocol}
Predicted and gold mentions are aligned by exact character offsets.
Following \cite{genest2025owner}, an unmatched gold mention (a false negative) and an unmatched predicted mention (a false positive) are each assigned a distinct sentinel label before computing AMI, so that detection errors are penalised rather than ignored.
We distinguish two protocols that must not be mixed in the same comparison.
In the \emph{typing-on-gold} protocol, the entity-typing component is applied to gold mention spans, which measures typing quality in isolation under perfect detection.
In the \emph{end-to-end} protocol, mentions are produced by the detector and the sentinel labels above are active, which is the setting comparable to the baselines.
All runs use a single fixed random seed ($42$); we note this as a limitation, since we do not report variance over seeds.

\subsection{Implementation Details}
\label{sec:exp:impl}
All experiments run on a single consumer-grade GPU with $6$\,GB of VRAM (CUDA~12.1, PyTorch~2.5.1 \cite{paszke2019pytorch}).
We use the HuggingFace Transformers library \cite{wolf2020transformers} for the baselines, sentence-transformers \cite{reimers2019sbert} for the Matryoshka embedder \cite{nussbaum2024nomic, kusupati2022matryoshka}, and scikit-learn \cite{pedregosa2011scikit} for the linear classifier, whose solver is LIBLINEAR \cite{fan2008liblinear}.
Latency is measured with two warmup passes followed by timed passes with GPU synchronisation, reporting p50, p95 and p99 per sentence.
Energy is measured with CodeCarbon \cite{lacoste2019codecarbon} in offline mode using the French grid carbon intensity; when hardware power counters are unavailable, we fall back to a thermal-design-power estimate and flag the measurement method in the report.
To keep the energy comparison fair, all systems are measured on the same machine within a single session, under comparable thermal conditions.

\subsection{Main Results}
\label{sec:exp:results}
% PLACEHOLDER TABLES. Toutes les cellules numeriques sont "--" en attendant les
% runs finaux. A remplir + mettre en gras le meilleur de chaque ligne/colonne.
% Source des chiffres a venir : outputs/results/ (aggregate_results.py).
\emph{[Narrative to be written once the numbers are in; the tables below give the target structure.]}

\begin{table*}[t]
\centering
\caption{Open-world NER. Adjusted Mutual Information (AMI, $\uparrow$) per dataset, end-to-end protocol. Dataset abbreviations follow \autoref{tab:datasets}; best per column in bold.}
\label{tab:main}
\footnotesize
\setlength{\tabcolsep}{4pt}
\begin{tabular}{l*{13}{c}c}
\toprule
Model & AI & Lit & Mus & Pol & Sci & WNUT & Rest & Movie & Fab & Bio & CoNLL & GUM & GEN & Avg \\
\midrule
OWNER    & -- & -- & -- & -- & -- & -- & -- & -- & -- & -- & -- & -- & -- & -- \\
GLiNER-S & -- & -- & -- & -- & -- & -- & -- & -- & -- & -- & -- & -- & -- & -- \\
GLiNER-M & -- & -- & -- & -- & -- & -- & -- & -- & -- & -- & -- & -- & -- & -- \\
GLiNER-L & -- & -- & -- & -- & -- & -- & -- & -- & -- & -- & -- & -- & -- & -- \\
GNER     & -- & -- & -- & -- & -- & -- & -- & -- & -- & -- & -- & -- & -- & -- \\
Opener   & -- & -- & -- & -- & -- & -- & -- & -- & -- & -- & -- & -- & -- & -- \\
\bottomrule
\end{tabular}
\end{table*}

\begin{table}[t]
\centering
\caption{Efficiency. Median per-sentence latency, energy and carbon per run, with parameter count. Lower is better for all but params.}
\label{tab:efficiency}
\small
\begin{tabular}{lrrrr}
\toprule
Model & Params & p50 (ms) & Energy (kWh) & CO\textsubscript{2} (g) \\
\midrule
OWNER     & -- & -- & -- & -- \\
GLiNER-S  & -- & -- & -- & -- \\
GLiNER-M  & -- & -- & -- & -- \\
GLiNER-L  & -- & -- & -- & -- \\
GNER      & -- & -- & -- & -- \\
Opener    & -- & -- & -- & -- \\
\bottomrule
\end{tabular}
\end{table}

\begin{table}[t]
\centering
\caption{Per-subdomain CrossNER (AMI, $\uparrow$). Best per column in bold.}
\label{tab:crossner}
\small
\begin{tabular}{lcccccc}
\toprule
Model & AI & Lit. & Music & Pol. & Sci. & Mean \\
\midrule
OWNER     & -- & -- & -- & -- & -- & -- \\
GLiNER-L  & -- & -- & -- & -- & -- & -- \\
GNER      & -- & -- & -- & -- & -- & -- \\
Opener    & -- & -- & -- & -- & -- & -- \\
\bottomrule
\end{tabular}
\end{table}

\subsection{Ablations}
\label{sec:exp:ablations}
% PLACEHOLDER. Table 4 a remplir : effet du contrastive (on/off), du class_weight
% balanced (on/off), typing-on-gold vs end-to-end, sweep dimension Matryoshka
% (64/128/256/512/768), choix du classifieur (GMM/LogReg/SVM).
\emph{[Ablation study to be added: contrastive on/off, balanced on/off, typing-on-gold vs end-to-end, Matryoshka dimension sweep, and classifier choice.]}
